% ==============================================================================
% Chapter 1: Introduction and Problem Formulation
% Owner: Member 1 (TV1) | Target length: ~4-5 pages
% ==============================================================================

\chapter{Giới thiệu bài toán và định hướng nghiên cứu}
\label{chap:intro}

% ------------------------------------------------------------------------------
\section{Bối cảnh và động lực nghiên cứu}
\label{sec:intro_context}

Các mô hình ngôn ngữ lớn (LLM) và mô hình ngôn ngữ thị giác lớn (LVLM, MLLM) hiện nay có thể xử lý tốt văn bản cũng như kết hợp văn bản với hình ảnh trong nhiều tác vụ hỏi đáp. Tuy nhiên, khi chỉ dựa vào tri thức tham số (\textit{parametric knowledge}) được lưu cố định trong bộ trọng số sau huấn luyện, các mô hình này vẫn gặp ba hạn chế cơ bản theo tổng hợp của Abootorabi và cộng sự~\cite{abootorabi2025ask}:
\begin{itemize}
    \item \textbf{Sinh thông tin sai sự thật (\textit{Hallucination}):} Mô hình có thể tạo ra câu trả lời trôi chảy về mặt diễn đạt nhưng sai lệch nội dung thực tế, nhất là với các thực thể hoặc chủ đề ít xuất hiện trong dữ liệu huấn luyện.
    \item \textbf{Tri thức lỗi thời (\textit{Outdated Knowledge}):} Kiến thức của mô hình dừng lại ở thời điểm khóa tập dữ liệu huấn luyện nên không phản ánh được thông tin mới cập nhật hay dữ liệu nội bộ của từng tổ chức.
    \item \textbf{Khó đối chiếu nguồn dẫn chứng (\textit{Lack of Verifiable Reasoning}):} Câu trả lời sinh trực tiếp từ trọng số mô hình không gắn kèm tài liệu tham chiếu cụ thể, gây khó khăn khi cần kiểm tra lại độ chính xác trong các lĩnh vực như y tế, tài chính hoặc kỹ thuật.
\end{itemize}

Phương pháp tạo sinh tăng cường truy xuất (RAG) giải quyết các hạn chế trên bằng cách kết hợp mô hình sinh với một kho dữ liệu bên ngoài. Trước khi trả lời, hệ thống tìm các tài liệu liên quan đến câu hỏi rồi đưa vào ngữ cảnh đầu vào của mô hình~\cite{abootorabi2025ask}. Cách làm này cho phép cập nhật kho dữ liệu độc lập mà không phải huấn luyện lại mô hình, đồng thời cung cấp tài liệu gốc để người dùng kiểm tra lại câu trả lời.

Dù vậy, phần lớn hệ thống RAG truyền thống chỉ làm việc trên văn bản thuần túy. Trong thực tế, dữ liệu lưu trữ thường nằm ở nhiều phương thức khác nhau như tài liệu PDF chứa bảng biểu và đồ thị, ảnh chụp, bảng dữ liệu có cấu trúc, âm thanh hoặc video~\cite{abootorabi2025ask,yeo2026universalrag}. Nếu chuyển toàn bộ các dạng dữ liệu này sang văn bản thông qua nhận dạng ký tự quang học (OCR) hoặc sinh chú thích ảnh tự động, hệ thống sẽ tốn nhiều thời gian tiền xử lý và làm mất thông tin về bố cục trang, vị trí tương đối trên biểu đồ hoặc diễn biến theo thời gian trong video~\cite{faysse2025colpali,yu2025visrag}.

Từ nhu cầu đó, hướng nghiên cứu \textbf{Multimodal RAG} (RAG đa phương thức) mở rộng quy trình truy xuất và tạo sinh sang nhiều loại dữ liệu khác nhau. Bài toán này không chỉ dừng ở việc đưa thêm ảnh hay video vào chỉ mục tìm kiếm, mà cần giải quyết đồng bộ từ bước chọn đúng kho dữ liệu, biểu diễn tài liệu giàu yếu tố thị giác, đến bước kiểm tra mâu thuẫn giữa các nguồn thông tin trước khi sinh câu trả lời~\cite{abootorabi2025ask,yeo2026universalrag,tian2025coremmrag}.

% ------------------------------------------------------------------------------
\section{Phát biểu bài toán Multimodal RAG}
\label{sec:intro_problem}

\subsection{Mô hình hình thức hóa}

Theo bài khảo sát của Abootorabi và cộng sự~\cite{abootorabi2025ask}, hệ thống Multimodal RAG nhận đầu vào là một truy vấn $q$ (dạng văn bản hoặc cặp văn bản kèm ảnh) cùng kho dữ liệu bên ngoài $D = \{d_1, d_2, \dots, d_n\}$ để sinh ra câu trả lời $r$. Mỗi phần tử dữ liệu $d_i \in D$ thuộc một phương thức $M_{d_i}$ xác định (văn bản, bảng biểu, hình ảnh, âm thanh hoặc video). Quy trình xử lý cơ bản gồm ba bước:

\begin{enumerate}
    \item \textbf{Mã hóa dữ liệu (\textit{Encoding}):} Mỗi đơn vị dữ liệu $d_i$ được đưa qua bộ mã hóa tương ứng với phương thức $M_{d_i}$ để thu được biểu diễn vector $z_i$:
    \begin{equation}
        z_i = \mathrm{Enc}_{M_{d_i}}(d_i), \qquad Z = \{z_1, z_2, \dots, z_n\}
        \label{eq:ch1_encode}
    \end{equation}
    Đối với tài liệu chứa nhiều phương thức trên cùng một trang (ví dụ trang PDF bài báo có cả chữ, bảng và hình vẽ), hệ thống có thể tách riêng từng thành phần để mã hóa hoặc chụp toàn bộ trang thành một ảnh rồi đưa trực tiếp vào mô hình ngôn ngữ thị giác~\cite{abootorabi2025ask,faysse2025colpali,yu2025visrag}.

    \item \textbf{Truy xuất ngữ cảnh (\textit{Retrieval}):} Bộ truy xuất tính độ tương đồng $s(e_q, z_i)$ giữa vector truy vấn $e_q$ và từng biểu diễn $z_i \in Z$, từ đó lấy ra tập ngữ cảnh $X \subseteq D$ vượt ngưỡng $\tau_{M_{d_i}}$ (hoặc chọn top-$k$ phần tử có điểm cao nhất):
    \begin{equation}
        X = \left\{ d_i \in D \;\middle|\; s(e_q, z_i) \ge \tau_{M_{d_i}} \right\}
        \label{eq:ch1_retrieve}
    \end{equation}
    Khi dữ liệu được chia thành nhiều kho riêng theo phương thức và mức độ chi tiết $\{\mathcal{C}_m\}$, Yeo và cộng sự~\cite{yeo2026universalrag} bổ sung một bộ định tuyến ở đầu quy trình để chọn tập kho phù hợp trước, sau đó mới chạy bộ truy xuất $\mathcal{T}_m(q; \mathcal{C}_m)$ trên từng kho đã chọn.

    \item \textbf{Tạo sinh câu trả lời (\textit{Generation}):} Mô hình tạo sinh $G$ (thường là LVLM hoặc MLLM) nhận truy vấn gốc $q$ cùng tập chứng cứ $X$ để sinh câu trả lời cuối:
    \begin{equation}
        r = G(q, X)
        \label{eq:ch1_generate}
    \end{equation}
\end{enumerate}

\subsection{Đặc tả đầu vào, đầu ra và ví dụ minh họa}

Từ mô hình trên, đầu vào và đầu ra của bài toán được xác định cụ thể như sau:
\begin{itemize}
    \item \textbf{Đầu vào:} (1) Truy vấn $q$ của người dùng dưới dạng văn bản hoặc văn bản kết hợp hình ảnh; (2) Kho dữ liệu bên ngoài $D$ gồm một hoặc nhiều kho thành phần (đoạn văn ngắn, tài liệu dài, bảng số liệu, ảnh trang PDF, ảnh tự nhiên, đoạn clip ngắn và video đầy đủ).
    \item \textbf{Đầu ra:} Câu trả lời $r$ bằng ngôn ngữ tự nhiên dựa trên bằng chứng truy xuất được (có thể kèm chỉ dẫn nguồn hoặc vùng ảnh liên quan), hoặc trả lời trực tiếp từ kiến thức sẵn có của mô hình nếu câu hỏi không cần tra cứu ngoài.
\end{itemize}

Bốn tình huống dưới đây cho thấy vì sao một kỹ thuật đơn lẻ khó đáp ứng được toàn bộ yêu cầu của bài toán:
\begin{itemize}
    \item \textbf{Chọn đúng phương thức và mức độ chi tiết cần tìm:} Với câu hỏi phổ thông như \textit{``Thủ đô của Pháp là gì?''}, mô hình có thể trả lời ngay mà không cần truy xuất. Câu hỏi về một tình huống ghi bàn chỉ cần đoạn clip ngắn vài phút thay vì cả video trận đấu dài hai giờ. Trong khi đó, câu hỏi \textit{``Có bao nhiêu tượng người trên đài tưởng niệm Michigan Soldiers and Sailors?''} lại cần kết hợp cả đoạn văn mô tả lẫn ảnh chụp vì mỗi nguồn riêng lẻ đều bị thiếu thông tin hoặc che khuất góc nhìn~\cite{yeo2026universalrag}.
    \item \textbf{Đọc hiểu tài liệu PDF có bố cục phức tạp:} Trên các trang infographic hoặc báo cáo tài chính dùng chữ cách điệu và bảng nhiều cột, công cụ OCR dễ nhận dạng sai ký tự (ví dụ đọc nhầm \textit{``Club Jetty''} thành \texttt{CLUBIETTY}) hoặc làm xáo trộn thứ tự hàng--cột. Giữ nguyên trang tài liệu dưới dạng ảnh giúp bảo toàn cấu trúc trình bày cho cả bước tìm kiếm lẫn bước đọc hiểu~\cite{yu2025visrag,faysse2025colpali}.
    \item \textbf{Đối khớp từ khóa truy vấn với chi tiết trên biểu đồ:} Khi câu hỏi yêu cầu tìm khung giờ tiêu thụ điện cao nhất trên biểu đồ năng lượng, bộ truy xuất cần liên kết từ \textit{``hour''} trong câu hỏi với cả phần chú giải lẫn trục ngang của đồ thị ngay trên ảnh trang~\cite{faysse2025colpali}.
    \item \textbf{Xử lý mâu thuẫn giữa các nguồn thông tin:} Trong hỏi đáp dựa trên ảnh, đoạn văn truy xuất được có thể mô tả một đối tượng khác có tên gần giống đối tượng trong ảnh, hoặc thông tin mới truy xuất được mâu thuẫn với kiến thức cũ trong trọng số mô hình. Khi đó, hệ thống cần bước kiểm tra và chọn lọc nguồn tin cậy thay vì ghép toàn bộ ngữ cảnh vào prompt một cách thụ động~\cite{tian2025coremmrag}.
\end{itemize}

% ------------------------------------------------------------------------------
\section{Các thách thức cốt lõi và bài toán con cần giải quyết}
\label{sec:intro_challenges}

Khi triển khai RAG trên dữ liệu đa phương thức, bốn bài toán con sau ảnh hưởng lớn nhất đến độ chính xác và chi phí vận hành của hệ thống:

\begin{enumerate}
    \item \textbf{Khoảng cách phương thức (\textit{Modality Gap}) và lựa chọn mức độ chi tiết (\textit{Granularity}):} Khi đưa văn bản, hình ảnh và video vào chung một không gian vector đa phương thức, các vector biểu diễn vẫn có xu hướng gom cụm theo từng loại phương thức thay vì trộn lẫn hoàn toàn theo nghĩa~\cite{liang2022mind,yeo2026universalrag}. Do truy vấn thường ở dạng văn bản, vector truy vấn nằm gần cụm văn bản hơn nên bộ tìm kiếm dễ bỏ sót ảnh và video liên quan. Ngoài ra, ngay trong cùng một phương thức, việc cố định một kích thước đơn vị lưu trữ cũng gây hạn chế: chia quá nhỏ làm thiếu ngữ cảnh cho câu hỏi suy luận nhiều bước, còn để quá dài lại đưa thêm nhiều nội dung nhiễu vào mô hình sinh~\cite{yeo2026universalrag}.

    \item \textbf{Chi phí tiền xử lý tài liệu và mất mát thông tin thị giác:} Quy trình bóc tách PDF truyền thống (phân tích bố cục, chạy OCR, tách bảng, sinh mô tả hình vẽ) gồm nhiều khâu nối tiếp, tốn thời gian lập chỉ mục và không thể tối ưu chung đầu--cuối. Sai sót ở bất kỳ khâu trung gian nào cũng truyền thẳng sang bước truy xuất và bước sinh câu trả lời~\cite{faysse2025colpali,yu2025visrag}.

    \item \textbf{Mâu thuẫn tri thức giữa các nguồn và thiên lệch phương thức:} Đầu ra của MLLM trong Multimodal RAG chịu ảnh hưởng cùng lúc từ hình ảnh đầu vào, tài liệu truy xuất bên ngoài và tri thức nội tại của mô hình. Hai dạng xung đột thường gặp là mâu thuẫn giữa tri thức nội tại với tài liệu truy xuất (PRKI) và mâu thuẫn giữa nội dung ảnh với văn bản truy xuất (VTKI). Nếu chỉ nối toàn bộ tài liệu tìm được vào đầu vào, mô hình rất dễ tin theo đoạn văn nhiễu hoặc trả lời theo thói quen có sẵn trong tham số~\cite{tian2025coremmrag,abootorabi2025ask}.

    \item \textbf{Đánh đổi giữa độ chi tiết biểu diễn, bộ nhớ chỉ mục và độ trễ:} Lưu mỗi trang tài liệu bằng hàng nghìn vector cấp vùng ảnh (\textit{patch}) giúp giữ được chi tiết chữ nhỏ và biểu đồ, nhưng làm dung lượng chỉ mục tăng lên hàng chục lần so với cách dùng một vector duy nhất cho cả trang~\cite{faysse2025colpali,yu2025visrag}. Ngược lại, nén cả trang thành một vector tiết kiệm bộ nhớ nhưng khó giữ hết chi tiết cục bộ; đồng thời việc ghép quá nhiều ảnh trang hoặc video dài vào bộ sinh cũng làm tăng thời gian suy luận và giảm độ chính xác do nhiễu~\cite{abootorabi2025ask,yu2025visrag}.
\end{enumerate}

% ------------------------------------------------------------------------------
\section{Chiến lược lựa chọn bài báo và cấu trúc báo cáo}
\label{sec:intro_strategy}

Để phân tích bài toán một cách có hệ thống thay vì tóm tắt rời rạc từng công trình, nhóm lựa chọn năm bài báo công bố tại ACL 2025, ICLR 2025 và ACL 2026 dựa trên \textbf{hai tiêu chí kết hợp}:

\begin{itemize}
    \item \textbf{Phân chia theo ba tầng nối tiếp của quy trình xử lý:} Nhóm sử dụng bài khảo sát \textit{Ask in Any Modality}~\cite{abootorabi2025ask} làm khung phân loại chung, sau đó khảo sát sâu ba chặng xử lý kế tiếp nhau trong một hệ thống Multimodal RAG: (i) \textit{UniversalRAG}~\cite{yeo2026universalrag} ở tầng trước truy xuất (\textit{Pre-Retrieval}), giải quyết việc chọn phương thức và mức độ chi tiết của kho dữ liệu; (ii) \textit{ColPali}~\cite{faysse2025colpali} và \textit{VisRAG}~\cite{yu2025visrag} ở tầng truy xuất và đọc hiểu tài liệu bằng hình ảnh (\textit{In-Corpus Visual Document RAG}); và (iii) \textit{CoRe-MMRAG}~\cite{tian2025coremmrag} ở tầng sau truy xuất (\textit{Post-Retrieval}), tập trung phát hiện và xử lý mâu thuẫn tri thức giữa các nguồn.
    \item \textbf{Đối chứng trực tiếp hai hướng thiết kế trên cùng một bài toán con:} Tại tầng tài liệu trực quan, nhóm đặt \textit{ColPali} (biểu diễn đa vector kết hợp tương tác muộn MaxSim) bên cạnh \textit{VisRAG} (biểu diễn đơn vector kết hợp mô hình sinh đọc nhiều ảnh trang) để so sánh rõ sự đánh đổi giữa độ chính xác tìm kiếm, dung lượng lưu trữ, tốc độ truy vấn và khả năng minh giải.
\end{itemize}

Bảng~\ref{tab:selected_papers} tóm tắt thông tin xuất bản và vai trò của năm bài báo trong báo cáo; Hình~\ref{fig:pipeline_overview} thể hiện mối liên hệ giữa năm công trình trong quy trình xử lý chung.

\begin{table}[H]
\centering
\small
\caption[Danh sách 5 bài báo trọng tâm và vai trò trong báo cáo]{Danh sách 5 bài báo trọng tâm và vai trò trong cấu trúc báo cáo.}
\label{tab:selected_papers}
\renewcommand{\arraystretch}{1.24}
\begin{tabularx}{\textwidth}{@{}M{0.8cm}P{3.5cm}P{2.5cm}P{2.4cm}L@{}}
\toprule
\rowcolor{tblhead}
\thead{STT} & \thead{Bài báo} & \thead{Tác giả} & \thead{Hội nghị / Năm} & \thead{Vai trò trong báo cáo} \\
\midrule
1 & \textit{Ask in Any Modality}~\cite{abootorabi2025ask} & Abootorabi et al. & ACL 2025 Findings & Khung phân loại 5 thành phần, phát biểu bài toán và tổng hợp bộ dữ liệu, độ đo. \\
\rowcolor{zebrabg}
2 & \textit{UniversalRAG}~\cite{yeo2026universalrag} & Yeo et al. & ACL 2026 (Oral) & Tầng Pre-Retrieval: định tuyến theo phương thức và mức độ chi tiết, khắc phục \textit{Modality Gap}. \\
3 & \textit{ColPali}~\cite{faysse2025colpali} & Faysse et al. & ICLR 2025 & Tầng Visual Doc Retrieval: biểu diễn đa vector, tương tác muộn MaxSim và bộ chuẩn ViDoRe. \\
\rowcolor{zebrabg}
4 & \textit{VisRAG}~\cite{yu2025visrag} & Yu et al. & ICLR 2025 & Tầng Visual Doc RAG: biểu diễn đơn vector gọn nhẹ và quy trình tạo sinh trực quan đầu--cuối. \\
5 & \textit{CoRe-MMRAG}~\cite{tian2025coremmrag} & Tian et al. & ACL 2025 & Tầng Post-Retrieval: đối chiếu và xử lý mâu thuẫn tri thức đa nguồn (PRKI và VTKI). \\
\bottomrule
\end{tabularx}
\end{table}

% Include overview pipeline diagram
% ==============================================================================
% Publication-grade TikZ Pipeline Diagram
% Guarantees aligned horizontal flow, non-overlapping inner cards, and centered links
% ==============================================================================
\begin{figure}[H]
\centering
\resizebox{\textwidth}{!}{%
\begin{tikzpicture}[
    >=Stealth,
    stagebox/.style={
        rectangle, rounded corners=4pt,
        draw=#1, fill=#1!8!white,
        line width=1.0pt, align=center,
        text width=3.45cm, inner sep=7pt,
        minimum height=5.05cm,
        drop shadow={shadow xshift=0.6pt, shadow yshift=-0.8pt, opacity=0.14}
    },
    endpoint/.style={
        rectangle, rounded corners=4pt,
        draw=#1, fill=#1!10!white,
        line width=1.0pt, align=center,
        text width=2.25cm, inner sep=6pt,
        minimum height=1.85cm,
        drop shadow={shadow xshift=0.6pt, shadow yshift=-0.8pt, opacity=0.14}
    },
    subcard/.style={
        rectangle, rounded corners=3pt,
        draw=tealacc!80!black, fill=white,
        line width=0.75pt, align=center,
        text width=4.0cm, inner sep=5pt,
        minimum height=1.35cm
    },
    flowarrow/.style={->, line width=1.25pt, color=primary!85},
    dashlink/.style={dashed, line width=0.85pt, color=muted}
]

% Stage 0: Input Query (aligned along y = 0)
\node[endpoint=muted] (query) at (0, 0) {
    {\small\bfseries\color{primary}Truy vấn}\\[4pt]
    {\footnotesize Văn bản /\\ Đa phương thức}
};

% Stage 1: Pre-Retrieval (UniversalRAG)
\node[stagebox=accent, right=0.7cm of query] (router) {
    {\footnotesize\bfseries\color{accent!90!black}TẦNG 1: PRE-RETRIEVAL}\\[6pt]
    {\small\bfseries\color{primary}UniversalRAG} {\scriptsize(ACL '26)}\\[6pt]
    {\footnotesize Multi-hot Router\\
    6 kho độc lập\\
    \textit{(Modality $\times$ Granularity)}}
};

% Stage 2: Vision-Space Document RAG container (exact same height & midline y=0)
\node[stagebox=tealacc, text width=4.75cm, right=0.75cm of router] (vision) {};
\node[font=\footnotesize\bfseries\color{tealacc!90!black}, anchor=north]
    at ($(vision.north)+(0,-0.22cm)$) {TẦNG 2: VISUAL DOC RAG};
\node[subcard, anchor=north] at ($(vision.north)+(0,-0.82cm)$) (colpali) {
    {\small\bfseries\color{primary}ColPali} {\scriptsize(ICLR '25)}\\[2pt]
    {\scriptsize Đa vector (1030$\times$128) + MaxSim}
};
\node[subcard, below=0.28cm of colpali] (visrag) {
    {\small\bfseries\color{primary}VisRAG} {\scriptsize(ICLR '25)}\\[2pt]
    {\scriptsize Đơn vector (4.5\,KB) + Visual Gen}
};

% Stage 3: Post-Retrieval (CoRe-MMRAG)
\node[stagebox=goldacc, right=0.75cm of vision] (core) {
    {\footnotesize\bfseries\color{goldacc!90!black}TẦNG 3: POST-RETRIEVAL}\\[6pt]
    {\small\bfseries\color{primary}CoRe-MMRAG} {\scriptsize(ACL '25)}\\[6pt]
    {\footnotesize Hòa giải xung đột 4 bước\\
    Kiểm chứng PRKI \& VTKI}
};

% Output Answer
\node[endpoint=tealacc!85!black, right=0.7cm of core] (ans) {
    {\small\bfseries\color{primary}Câu trả lời}\\[4pt]
    {\footnotesize Trung thực \&\\ Có căn cứ}
};

% Top Survey Taxonomy banner spanning exact full width with generous vertical padding
\coordinate (bannerTopLeft) at ($(query.west |- router.north)+(0, 2.00cm)$);
\coordinate (bannerBotRight) at ($(ans.east |- router.north)+(0, 0.62cm)$);
\node[
    rectangle, rounded corners=4pt,
    draw=primary, fill=primary,
    line width=1.0pt,
    fit=(bannerTopLeft)(bannerBotRight),
    inner sep=0pt,
    drop shadow={shadow xshift=0.6pt, shadow yshift=-0.8pt, opacity=0.15}
] (survey) {};
\node[align=center, text=white] at (survey.center) {
    {\small\bfseries Khung phân loại tổng quan (Taxonomy 5 thành phần): \textit{Ask in Any Modality} (ACL 2025 Findings)}\\[3pt]
    {\footnotesize\color{white!90!primary} Retrieval $\;\bullet\;$ Fusion $\;\bullet\;$ Augmentation $\;\bullet\;$ Generation $\;\bullet\;$ Training}
};

% Strictly horizontal main pipeline arrows
\draw[flowarrow] (query.east) -- (router.west);
\draw[flowarrow] (router.east) -- (vision.west);
\draw[flowarrow] (vision.east) -- (core.west);
\draw[flowarrow] (core.east) -- (ans.west);

% Dead-center vertical dashed mapping lines from survey banner to each stage
\draw[dashlink, ->] (router.north |- survey.south) -- (router.north);
\draw[dashlink, ->] (vision.north |- survey.south) -- (vision.north);
\draw[dashlink, ->] (core.north |- survey.south) -- (core.north);

\end{tikzpicture}%
}
\caption[Định vị 5 bài báo nghiên cứu trong đường ống Multimodal RAG]{Định vị 5 bài báo nghiên cứu trong đường ống xử lý Multimodal RAG đầu--cuối.}
\label{fig:pipeline_overview}
\end{figure}


Nội dung tiếp theo của báo cáo gồm năm chương. Chương~\ref{chap:background} trình bày hệ thống phân loại của lĩnh vực, các hướng tiếp cận cơ sở và vị trí của bốn bài báo kỹ thuật trên khung phân loại đó. Chương~\ref{chap:methods} phân tích chi tiết kiến trúc, công thức mô hình hóa và khả năng kết hợp của bốn phương pháp theo ba tầng xử lý. Chương~\ref{chap:experiments} tổng hợp hệ thống bộ dữ liệu, độ đo, xét điều kiện so sánh giữa các kết quả thực nghiệm và đánh giá phản biện từng phương pháp. Chương~\ref{chap:reproduction} trình bày phần thực nghiệm kiểm chứng bổ sung, và Chương~\ref{chap:conclusion} tóm tắt các kết luận chính cùng những hướng nghiên cứu mở.
